\scanpage{139}
\deleted{\unsure{du fait de l'homogénéité|à part d'homogénéité} \unsure{il vient,|évidente,} qui permet}

Remplaçons dans cette formule $A$ par $\lambda A$, $B$ par $\mu B$ ($\lambda$ et $\mu$ nombres $\unsure{>| \geq}0$) on obtient :

\[
\deleted{\unsure{\lambda^p\mu^q=\frac{\lambda^p\mu^q}{r}S_r(AB)}}
\inserted{\frac{(\lambda\mu)^r}{r}S_r(AB)}
\leq\frac{\lambda^p}{p}S_p(A)+\frac{\mu^q}{q}S_q(B)
\]

\deleted{Exprimons que cette inégalité doit être vérifiée quels que soient $x,y\geq 0$.}

% slanted
\margin{\unsure{Ici lemme que} la condition du \unsure{n°} implique $c^{1/r}\leq a^{1/p}b^{1/q}$ i.e \uline{$c\leq a^\alpha b^\beta$}}

Pour simplifier, posons $a=S_p(A)$, $b=S_q(B)$, $c=S_{\unsure{r|p}}(C)$, on aura

\[
c\leq\mathop{\mathrm{Inf}}\limits_{\substack{\lambda\geq0\\\mu\geq0}}
\frac{a\frac{\lambda^p}{p}+b\frac{\mu^q}{q}}
{\deleted{\unsure{\frac{\lambda^p\mu^q}{r}}}\frac{(\lambda\mu)^r}{r}}
\]

\deleted{Posons $\unsure{x^p=\lambda,y^q}$}\deleted{$x=\lambda^p,\ y=\mu^q$,}

Posons $x=\lambda^p$, $y=\mu^q$, $\alpha=\frac rp$, $\beta=\frac rq$ d'où \unsure{,} \uline{$\frac rp+\frac rq=1$},

On a

\[
c\leq\mathop{\mathrm{Inf}}\limits_{\substack{x\unsure{>|\geq}0\\y>0}}
\frac{\alpha a\deleted{\unsure{\frac rp}}x+\deleted{\unsure{\frac rq}}\beta by}{x^\alpha y^\beta}
=\mathop{\mathrm{Inf}}\limits_{\substack{x\geq0\\y\geq0\\x^\alpha y^\beta=1}}\alpha a x+\beta by
\]

Or on vérifie aussitôt que la \deleted{minimum}\inserted{borne inférieure} du second membre est atteinte (car l'en\unsure{s.} des points $(x,y)$ tels que $x\geq0,y\geq0$ $ax+by\leq M$ -- valeur donnée -- est compact) elle est donc atteinte en un point $(x,y)$ tel que les équations \deleted{$a\,dx+b\,dy=0$}\inserted{$d(\alpha ax+\beta by)=0$}, $d(x^\alpha y^\beta)=0$ \deleted{soient compatibles} \inserted{aient une solution non nulle en} $dx,dy$  ;

or $d(x^\alpha y^\beta)=\alpha x^{\alpha-1}y^\beta\,dx+\beta y^{\beta-1}x^\alpha\,dy$, \deleted{i.e. tel que $\unsure{\det(\alpha a\quad\beta b\quad\alpha a}$} la condition d'extremum est donc $\frac{\alpha a}{\alpha x^{\alpha-1}y}=\frac{\beta b}{\beta x^\alpha y^{\beta-1}}$,

soit $ax=by$ \deleted{\unsure{$=$}}, \deleted{\unsure{et comme on a}} \inserted{d'où $ax=by=\frac{\alpha ax+\beta by}{\alpha+\beta}$} ; \deleted{$\unsure{ax+by=\frac{\alpha ax}{\alpha}=\frac{\alpha ax+\beta by}{\alpha}}$} Comme $\alpha+\beta=1$, le minimum cherché sera alors \deleted{\unsure{$ax+by$|$\alpha x+b$}} égal à $ax=by$ ; compte tenu de la condi%
% archive sheet number: 140
